\section{Grounded transport}\label{sec:gt}

\subsection{Ordinary relations and grounded witnesses}\label{sec:witnesses}
An ordinary relational judgement $R(x,y) : \Prop$ says that $x$ and $y$ are related. A grounded judgement $\GTrans_c(x,y) : \Type$ is indexed by a declared \emph{transport problem} $c$ and its inhabitants carry inspectable evidence of why $x$ may be transported to $y$. We make this precise for the class of problems that includes the running example.

\begin{defi}[Transport problem]\label{def:problem}
A \emph{transport problem} $c$ consists of
\begin{itemize}
\item a span $A \xleftarrow{\back} \Sigma \xrightarrow{\fwd} B$ with Boolean valuations $\Phi_A$, $\Phi_B$ and a seam valuation $\Phi_\Sigma : \Sigma \to \bool$, and observation maps $M_A : A \to O_A$, $M_B : B \to O_B$;
\item a \emph{declared lineage record} $\ell$ (a finite structure, \cref{sec:lineage});
\item a type $\mathrm{Mnt}$ of maintenance evidence, a type $\mathrm{Mnt}^\bot$ of located maintenance failures, and a function $\mathrm{Mnt}^\bot \to \mathrm{Mnt} \to \bot$ saying that a failure refutes any evidence;
\item the same triple $(\mathrm{Aut}, \mathrm{Aut}^\bot, \mathrm{Aut}^\bot \to \mathrm{Aut} \to \bot)$ for authority evidence.
\end{itemize}
\end{defi}
The generic record fixes no institutional record format: maintenance and authority evidence are abstract types supplied by the application. Write $\Adm(M,\Phi) :\equiv \exists h.\ \Phi = h \circ M$ (admissibility of $\Phi$ with respect to observation $M$).

\begin{defi}[Certificates]\label{def:certs}
For a problem $c$ we have the following records, all in $\Type$.
\begin{enumerate}
\item An \emph{exhibition certificate} $\mathsf{Exh}(c)$: a record identifier; an injective assignment $\iota : \Sigma \to \nat$ of stable identifiers; a declared process; a type $\mathrm{Ret}$ of retained records with $\mathsf{retain} : \Sigma \to \mathrm{Ret}$ and $\mathsf{eval} : \mathrm{Ret} \to \bool$ such that $\Phi_\Sigma = \mathsf{eval} \circ \mathsf{retain}$ (the seam valuation is a function of what the record retains); a declared coverage interval; and a custodian identifier.
\item A \emph{lineage certificate} $\mathsf{Lin}$: a lineage graph $g$, a coverage record with its declared gaps, and an issued verdict, as \emph{data}; and proofs that $g$ is well-formed and satisfies L1--L3 (\cref{sec:lineage}), and that the gaps are empty.
\item A \emph{construction certificate} $\mathsf{Con}(c)$: an element $\sigma_0 \in \Sigma$; an $e \in \mathsf{Exh}(c)$; a $\lambda \in \mathsf{Lin}$ whose graph \emph{is} $\ell$; and the \emph{grounding equations} $\Phi_A(\back(r)) = \Phi_\Sigma(r)$ and $\Phi_B(\fwd(r)) = \Phi_\Sigma(r)$ for all $r$.
\item A \emph{certified transport} $\mathsf{Cert}(c) :\equiv \bigl(\Adm(M_A,\Phi_A) \wedge \Adm(M_B,\Phi_B)\bigr) \times \mathsf{Con}(c) \times \mathrm{Mnt} \times \mathrm{Aut}$.
\end{enumerate}
\end{defi}

\begin{defi}[States and grounded witnesses]\label{def:gtrans}
The \emph{states} of $c$ are $S_c :\equiv A + B$, with $\val(\mathsf{inl}\,a) = \Phi_A(a)$ and $\val(\mathsf{inr}\,b) = \Phi_B(b)$. The \emph{grounded witnesses} $\GTrans_c(s,t)$ are generated by
\begin{gather*}
\frac{}{\mathsf{id}_s : \GTrans_c(s,s)}\qquad
\frac{v : \GTrans_c(s,t) \quad w : \GTrans_c(t,u)}{v \cdot w : \GTrans_c(s,u)}\\[4pt]
\frac{r \in \Sigma \quad \kappa \in \mathsf{Cert}(c)}{\mathsf{step}(r,\kappa) : \GTrans_c(\mathsf{inl}\,\back(r),\ \mathsf{inr}\,\fwd(r))}
\end{gather*}
The \emph{crossing witnesses} $\Cross_c(s,t)$ are generated by the same rules except that $\mathsf{step}(r)$ requires only the two grounding equations at $r$ and no certificate.
\end{defi}

A witness is thus a directed path of crossings, each carrying a full certified transport. States are the disjoint union of the two coordinates; a crossing at $r$ leads from the $A$-side image $\back(r)$ to the $B$-side image $\fwd(r)$.

\paragraph{What survives extraction.}
The data-bearing components are: the seam element and its inhabitant; the exhibition record identifier, seam identifiers, process, retained-record function and evaluator, coverage interval and custodian; the lineage graph, coverage record with its gaps, and issued verdict; and the maintenance and authority evidence (to the extent the application supplies them in $\Type$). The proofs that these records satisfy their specifications---injectivity, factorisation, well-formedness, L1--L3, the grounding equations---certify the data but are erased.

\subsection{Why assessment is problem-indexed}\label{sec:indexed}
Two things must be kept apart: formal reflexivity of the transport relation, and the assessment of a proposed concrete transport problem.
\begin{quote}
\emph{Relational reflexivity is a property of the formal rewriting structure. It is not a certificate for every concrete process whose endpoints happen to be equal.}
\end{quote}
The identity witness $\mathsf{id}_s$ needs no certificate, so a judgement ``$s$ may be transported to $s$'' is always available formally. Yet a proposed \emph{process} that leaves a state representationally unchanged may have altered its provenance, calibration or authority; a pair-indexed assessment $\Refu(x,x)$ would contradict formal reflexivity, whereas a problem-indexed assessment can reject a proposed concrete process from $x$ to $x$ without doing so. Assessments (\cref{sec:assessment}) are therefore indexed by a transport \emph{problem}, and pair-level witnesses are \emph{derived} from a certified assessment.

\subsection{Identity and composition}\label{sec:laws}
Abstracting what all instances share gives:
\begin{defi}[Transport structure]\label{def:structure}
A \emph{transport structure} $G$ for a relation $R$ on $X$ is a family $\GTrans_G : X \to X \to \Type$ with an identity $\mathsf{id}_x : \GTrans_G(x,x)$, a composition $\GTrans_G(x,y) \to \GTrans_G(y,z) \to \GTrans_G(x,z)$, and a soundness map $\GTrans_G(x,y) \to R(x,y)$.
\end{defi}
$\GTrans_c$ is a transport structure for $R_c(s,t) :\equiv \val(s) = \val(t)$ (soundness is \cref{thm:soundness}), and so is $\Cross_c$.

\begin{thm}[Grounded reflexivity and composition]\label{thm:refl}
For every transport structure, $\mathsf{id}_x : \GTrans_G(x,x)$ and $\GTrans_G(x,y) \to \GTrans_G(y,z) \to \GTrans_G(x,z)$. In particular both hold for $\GTrans_c$, with $\mathsf{id}_s$ and $v \cdot w$. \emph{(Coq: \coq{grounded_transport_refl}, \coq{grounded_transport_compose}; mechanised, immediate from the definition.)}
\end{thm}

\paragraph{Algebraic laws.}
$\GTrans_c$ is a \emph{free} structure, so identity and associativity do not hold by literal equality: $\mathsf{id} \cdot w$ is not syntactically $w$. We therefore state them modulo an explicit equivalence. Let $\mathsf{crossings}(w)$ be the list of pairs $(r,\kappa)$ occurring in $w$, defined by $\mathsf{crossings}(\mathsf{id}) = [\,]$, $\mathsf{crossings}(\mathsf{step}(r,\kappa)) = [(r,\kappa)]$, and $\mathsf{crossings}(v \cdot w) = \mathsf{crossings}(v) \mathbin{+\!\!+} \mathsf{crossings}(w)$. Define
\[
w_1 \weq w_2 \;:\Longleftrightarrow\; \mathsf{crossings}(w_1) = \mathsf{crossings}(w_2).
\]
This forgets bracketing and identity padding and \emph{nothing else}: every seam element and every certificate is compared, and the comparison is by (dependent) equality.

\begin{thm}[Witness algebra]\label{thm:algebra}
$\weq$ is an equivalence relation on $\GTrans_c(s,t)$; composition respects it ($v \weq v'$ and $w \weq w'$ imply $v \cdot w \weq v' \cdot w'$); and
\[
\mathsf{id} \cdot w \weq w, \qquad w \cdot \mathsf{id} \weq w, \qquad (w_1 \cdot w_2) \cdot w_3 \weq w_1 \cdot (w_2 \cdot w_3).
\]
Moreover $w_1 \weq w_2$ implies that $w_1$ and $w_2$ have the same extractable summary (\cref{sec:multiplicity}). \emph{(Coq: \coq{WEquiv_id_left}, \coq{WEquiv_id_right}, \coq{WEquiv_assoc}, \coq{WEquiv_compose}, \coq{WEquiv_summary}.)}
\end{thm}
\begin{proof}
$\weq$ is the kernel of the function $\mathsf{crossings}$, hence an equivalence. If $v \weq v'$ and $w \weq w'$ then
$\mathsf{crossings}(v \cdot w) = \mathsf{crossings}(v) \mathbin{+\!\!+} \mathsf{crossings}(w) = \mathsf{crossings}(v') \mathbin{+\!\!+} \mathsf{crossings}(w') = \mathsf{crossings}(v' \cdot w')$.
The three laws are the corresponding list laws: $[\,] \mathbin{+\!\!+} l = l$, $l \mathbin{+\!\!+} [\,] = l$, and associativity of $\mathbin{+\!\!+}$, applied to $\mathsf{crossings}$. Finally the summary of a witness is computed pointwise from $\mathsf{crossings}(w)$ (by induction on $w$, using $\mathsf{map}(f, l \mathbin{+\!\!+} l') = \mathsf{map}(f,l) \mathbin{+\!\!+} \mathsf{map}(f,l')$), so it depends only on that list.
\end{proof}
The equivalence is deliberately intensional: algebraically, $\weq$ is equality of morphisms in the free category generated by the certified crossings. No coarser semantic equivalence, identifying different certificates that produce the same external behaviour, is developed. Consequently the states of $S_c$, with $\weq$-classes of witnesses as morphisms, form a category \cite{maclane1998}. No quotient type is constructed in the development and we claim nothing stronger; in particular we do not claim a category of witnesses under literal equality. The same proof applies to $\Cross_c$.

\begin{rem}\label{rem:degenerate}
In a single seam, composition of non-identity crossing witnesses is degenerate: a crossing leads from the $A$-side to the $B$-side and has no continuation, so every crossing witness is an identity or a single crossing (\coq{opath_inv}). Composition earns its keep through contexts and temporal evolution (\cref{sec:contexts,sec:dynamic}), not through long chains inside one seam.
\end{rem}

\subsection{Contexts at two levels}\label{sec:contexts}
Ordinary $\Proper$ preserves a relation. A context that preserves the \emph{warrant} for the relation must map witnesses, and there are two distinct things it could map: the crossings that survive the first erasure, or the complete certified witnesses. Conflating them would overstate what a context preserves, so we treat both and relate them.

\paragraph{Generic witness-preserving contexts.}
\begin{defi}[Witness-preserving context]\label{def:gp}
Let $G_X$ and $G_Y$ be transport structures. A function $F : X \to Y$ is \emph{witness-preserving} if it comes with a map $\mathsf{map}_F : \GTrans_{G_X}(x,y) \to \GTrans_{G_Y}(Fx,Fy)$ satisfying
\[
\mathsf{map}_F(\mathsf{id}_x) = \mathsf{id}_{Fx}, \qquad \mathsf{map}_F(v \cdot w) = \mathsf{map}_F(v) \cdot \mathsf{map}_F(w).
\]
\end{defi}
The laws are equalities; for maps defined by structural recursion on witnesses they hold by unfolding. (Coq: \coq{GroundedProper}, a definition parametric in the two structures.) What such a context preserves depends entirely on the structures: at $G = \Cross_c$ it preserves crossings only; at $G = \GTrans_c$ it preserves complete certificates.

\begin{thm}[Composition]\label{thm:gpcomp}
The identity is witness-preserving, and if $F : X \to Y$ and $H : Y \to Z$ are, so is $H \circ F$, with $\mathsf{map}_{H \circ F} = \mathsf{map}_H \circ \mathsf{map}_F$. \emph{(Coq: \coq{gp_id}, \coq{gp_comp}, \coq{grounded_proper_compose}.)}
\end{thm}
\begin{proof}
For the identity take $\mathsf{map} = \mathrm{id}$. For $H \circ F$ put $\mathsf{map}_{H\circ F}(w) = \mathsf{map}_H(\mathsf{map}_F(w))$, well-typed since $\mathsf{map}_F(w) : \GTrans_{G_Y}(Fx,Fy)$. Identity: $\mathsf{map}_H(\mathsf{map}_F(\mathsf{id}_x)) = \mathsf{map}_H(\mathsf{id}_{Fx}) = \mathsf{id}_{HFx}$. Composition: $\mathsf{map}_H(\mathsf{map}_F(v \cdot w)) = \mathsf{map}_H(\mathsf{map}_F(v) \cdot \mathsf{map}_F(w)) = \mathsf{map}_H(\mathsf{map}_F(v)) \cdot \mathsf{map}_H(\mathsf{map}_F(w))$.
\end{proof}

\paragraph{Level 1: crossing-proper contexts.}
A witness-preserving map on crossing witnesses says nothing about certificates. To relate the two levels we also track the proof-independent content of a crossing witness, the list $\mathsf{cr}(p)$ of seam elements it crosses ($\mathsf{cr}(\mathsf{id}) = [\,]$, $\mathsf{cr}(\mathsf{step}(r)) = [r]$, $\mathsf{cr}(p \cdot q) = \mathsf{cr}(p) \mathbin{+\!\!+} \mathsf{cr}(q)$). Two crossing witnesses may differ in the proofs of their grounding equations; in Coq, which has no definitional proof irrelevance, comparing $\mathsf{cr}$ is the available proof-independent comparison. The grounding equations certify a crossing inside Coq but are erased on extraction; the seam-element sequence is the proof-independent crossing data retained by the extracted representation.
\begin{defi}[Crossing-proper]\label{def:crossingproper}
$F : S_c \to S_{c'}$ is \emph{crossing-proper} if it is witness-preserving between $\Cross_c$ and $\Cross_{c'}$ and comes with an \emph{action} $\alpha : \mathsf{List}(\Sigma) \to \mathsf{List}(\Sigma')$ on crossed seam elements such that $\mathsf{cr}(\mathsf{map}_F(p)) = \alpha(\mathsf{cr}(p))$ for all crossing witnesses $p$. (Coq: \coq{CrossingProper}.)
\end{defi}

\paragraph{Level 2: certificate-proper contexts.}
\begin{defi}[Certificate-proper]\label{def:certproper}
$F : S_c \to S_{c'}$ is \emph{certificate-proper} if it is crossing-proper, with action $\alpha$, and is witness-preserving between $\GTrans_c$ and $\GTrans_{c'}$ with a map $\mathsf{map}^{\mathrm{cert}}_F$ that \emph{commutes with the first erasure up to the action}:
\[
\mathsf{cr}\bigl(\erase_1(\mathsf{map}^{\mathrm{cert}}_F(w))\bigr) = \alpha\bigl(\mathsf{cr}(\erase_1(w))\bigr) \qquad\text{for all } w \in \GTrans_c(s,t).
\]
(Coq: \coq{CertificateProper}.)
\end{defi}
The commutation condition is the content of ``preserves certificates'': the certified witness produced by the context crosses exactly the seam elements that the crossing-level action prescribes, so the two levels describe the same transport.

\begin{thm}[Two levels of contextual preservation]\label{thm:twolevels}
(a) Identity is crossing-proper and certificate-proper; and if $F$ and $H$ are crossing-proper (resp.\ certificate-proper) then so is $H \circ F$. \emph{(Coq: \coq{crossing_proper_id}, \coq{crossing_proper_comp}, \coq{certificate_proper_id}, \coq{certificate_proper_comp}.)}
(b) Every certificate-proper context is crossing-proper, and both are proper for the erased relations:
\[
\text{certificate-proper} \;\Longrightarrow\; \text{crossing-proper} \;\Longrightarrow\; \Proper,
\]
that is, $\erased_{\GTrans_c}(s,t) \Rightarrow \erased_{\GTrans_{c'}}(Fs,Ft)$ and $\erased_{\Cross_c}(s,t) \Rightarrow \erased_{\Cross_{c'}}(Fs,Ft)$. \emph{(Coq: \coq{certificate_proper_erased}, \coq{crossing_proper_erased}, \coq{contextual_chain}.)}
(c) Nesting: for composable contexts $C_1,\dots,C_n$ of either level, the composite is of that level and its witness map is $\mathsf{map}_{C_n}\circ\cdots\circ\mathsf{map}_{C_1}$, so
$\mathsf{map}_{C_n}\cdots\mathsf{map}_{C_1}(w) : \GTrans\bigl(C_n[\cdots C_1[x]],\ C_n[\cdots C_1[y]]\bigr)$. \emph{(Coq: \coq{contextual_preservation}, \coq{contextual_preservation_nested}.)}
\end{thm}
\begin{proof}
(a) The witness maps compose by \cref{thm:gpcomp}; the action of a composite is $\alpha_H \circ \alpha_F$, and $\mathsf{cr}(\mathsf{map}_H(\mathsf{map}_F(p))) = \alpha_H(\mathsf{cr}(\mathsf{map}_F(p))) = \alpha_H(\alpha_F(\mathsf{cr}(p)))$. For certificate-proper contexts, using the commutation of $H$ and then of $F$,
$\mathsf{cr}(\erase_1(\mathsf{map}^{\mathrm{cert}}_H(\mathsf{map}^{\mathrm{cert}}_F(w)))) = \alpha_H(\mathsf{cr}(\erase_1(\mathsf{map}^{\mathrm{cert}}_F(w)))) = \alpha_H(\alpha_F(\mathsf{cr}(\erase_1(w))))$.
(b) The first implication is the projection to the crossing-proper part, which is part of \cref{def:certproper}. For $\Proper$: from $w \in \GTrans_c(s,t)$ obtain $\mathsf{map}^{\mathrm{cert}}_F(w) \in \GTrans_{c'}(Fs,Ft)$, hence $\erased_{\GTrans_{c'}}(Fs,Ft)$, and by erasure soundness $\val(Fs) = \val(Ft)$; likewise for crossing witnesses. (c) Induction on $n$, using (a).
\end{proof}
The converses are not claimed. A crossing-proper context need not lift to certificates (it has no certificate to lift), and an ordinary $\Proper$ instance need not yield any witness map: a context may preserve endpoint equivalence while discarding or fabricating lineage.

\paragraph{Evidence lifting.}
What extra structure yields certificate-proper contexts? Certificates (exhibition, lineage, coverage, maintenance, authority) must themselves be transported.
\begin{defi}[Problem morphism with evidence lift]\label{def:pm}
A \emph{problem morphism} $\varphi : c \to c'$ consists of maps $\varphi_\Sigma$, $\varphi_A$, $\varphi_B$ commuting with the projections ($\back' \circ \varphi_\Sigma = \varphi_A \circ \back$, $\fwd' \circ \varphi_\Sigma = \varphi_B \circ \fwd$); a proof that the grounding equations at $r$ imply those at $\varphi_\Sigma(r)$; and an \emph{evidence lift} $\mathsf{lift} : \mathsf{Cert}(c) \to \mathsf{Cert}(c')$. Its \emph{state map} is $\mathsf{inl}\,a \mapsto \mathsf{inl}(\varphi_A a)$, $\mathsf{inr}\,b \mapsto \mathsf{inr}(\varphi_B b)$. (Coq: \coq{ProblemMorphism}.)
\end{defi}
\begin{thm}[Evidence lifting gives certificate-proper contexts]\label{thm:lifting}
The state map of a problem morphism $\varphi : c \to c'$ is certificate-proper, with action $\mathsf{map}(\varphi_\Sigma)$. \emph{(Coq: \coq{pm_certificate_proper}, \coq{pm_crossing_proper}, \coq{pm_commute}.)}
\end{thm}
\begin{proof}
Define both witness maps by recursion. On crossing witnesses: $\mathsf{id} \mapsto \mathsf{id}$; $p \cdot q \mapsto \mathsf{map}(p) \cdot \mathsf{map}(q)$; and $\mathsf{step}(r) \mapsto \mathsf{step}(\varphi_\Sigma r)$, whose grounding equations come from the preservation hypothesis, transported along the two commuting squares so that its endpoints are the images of the endpoints of $\mathsf{step}(r)$. On certified witnesses the same, with $\mathsf{step}(r,\kappa) \mapsto \mathsf{step}(\varphi_\Sigma r, \mathsf{lift}(\kappa))$, cast along the same squares. The two identity and composition laws hold by definition, and $\mathsf{cr}$ of the image of a step is $[\varphi_\Sigma r] = \mathsf{map}(\varphi_\Sigma)([r])$ (casts do not change the crossed seam element), so the crossing action is $\mathsf{map}(\varphi_\Sigma)$ and, by induction on $w$ (the composition case is $\mathsf{map}_{\mathrm{list}}(f)(l \mathbin{+\!\!+} l') = \mathsf{map}_{\mathrm{list}}(f)(l) \mathbin{+\!\!+} \mathsf{map}_{\mathrm{list}}(f)(l')$), the commutation condition holds for certified witnesses.
\end{proof}
The evidence lift is the extra structure. It is the reason certificate-level preservation is a stronger claim than crossing-level preservation, and \cref{sec:dynamic} shows an evolution for which the lift exists and one for which it cannot.

\begin{rem}[What each level shows]\label{rem:crossinglevel}
Crossing-proper contexts preserve certified \emph{crossings} compositionally. Certificate-proper contexts additionally preserve complete evidential certificates, and require an evidence lift. The seam evolutions of \cref{sec:dynamic} are crossing-proper given preservation of grounding on the image; they are certificate-proper only under an evidence lift. This fits the paper's central lesson: preservation of structural or extensional transport does not automatically preserve warrant.
\end{rem}
The development provides semantic combinators and instances, not an automated witness-carrying rewriting tactic. A tactic consuming the kernel should fail only when no suitable witness was supplied or constructed, never because the semantics needed weakening.
